% RACER-GT user guide (English).
\documentclass[12pt,a4paper,oneside]{article}
% English setup. Compiled with XeLaTeX so that both languages share one engine
% and one font family, which keeps the two PDFs typographically comparable.
\usepackage{fontspec}
\setmainfont{Noto Serif}
\setsansfont{Noto Sans}
\setmonofont{Noto Sans Mono}

% Author-year citations. Loaded here rather than in _common because it must
% precede hyperref, and because the Chinese manuscript uses numeric citations
% that natbib's author-year mode would reject. Bibliographies are hand-written
% thebibliography blocks: biblatex/biber is absent from a BasicTeX install, and
% a document set that only builds on a full TeX Live is not reproducible enough
% for this project.
\usepackage[authoryear,round]{natbib}

% Single source of truth for version and date across every RACER-GT document.
% Regenerate nothing by hand: bump here, rebuild, and every PDF agrees.
\newcommand{\racerversion}{1.1.0}
\newcommand{\racerdate}{2026-07-27}
\newcommand{\racerauthor}{Yi-Hao Lai}
\newcommand{\raceraffil}{Department of Finance, Da-Yeh University}
\newcommand{\raceremail}{yhlai@mail.dyu.edu.tw}
\newcommand{\racerrepo}{https://github.com/yhlai0911/RACER-GT}

% Shared package set and mathematical macros for every RACER-GT document,
% in both languages. Language-specific font and theorem-name setup lives in
% _preamble-zh.tex and _preamble-en.tex.
%
% Font policy: the build depends only on the five Noto families below. They are
% installable on Linux CI (fonts-noto-core, fonts-noto-cjk) and on macOS via
% Homebrew casks. Do not introduce a font outside this set without adding it to
% .github/workflows/docs.yml as well, or the PDF build stops being reproducible.

\usepackage[margin=2.35cm,headheight=15pt]{geometry}
\usepackage{amsmath,amssymb,amsthm,mathtools,bm}
\usepackage{booktabs,longtable,tabularx,array,multirow}
\usepackage{graphicx,float,caption,subcaption}
\usepackage{enumitem}
\usepackage{xcolor}
\usepackage{microtype}
\usepackage{setspace}
\usepackage{fancyhdr}
\usepackage{hyperref}
\usepackage[nameinlink,noabbrev]{cleveref}
\usepackage{algorithm}
\usepackage{algpseudocode}
\usepackage{listings}
\usepackage{tikz}
\usetikzlibrary{arrows.meta,positioning,shapes.geometric,fit,calc}

\hypersetup{
  colorlinks=true,
  linkcolor=black,
  citecolor=black,
  urlcolor=blue,
  pdfauthor={\racerauthor}
}

\pagestyle{fancy}
\fancyhf{}
\rhead{v\racerversion}
\cfoot{\thepage}
\setstretch{1.28}
\setlength{\parskip}{0.35em}
\setlist[itemize]{leftmargin=2em,itemsep=0.25em,topsep=0.3em}
\setlist[enumerate]{leftmargin=2.2em,itemsep=0.25em,topsep=0.3em}

% ---------------------------------------------------------------- mathematics
\newcommand{\E}{\mathbb{E}}
\newcommand{\Var}{\operatorname{Var}}
\newcommand{\Cov}{\operatorname{Cov}}
\newcommand{\tr}{\operatorname{tr}}
\newcommand{\diag}{\operatorname{diag}}
\newcommand{\argmin}{\operatorname*{arg\,min}}
\newcommand{\argmax}{\operatorname*{arg\,max}}
\newcommand{\one}{\bm{1}}
\newcommand{\Rset}{\mathbb{R}}
\newcommand{\R}{\mathbb{R}}
\newcommand{\plim}{\operatorname*{plim}}
\newcommand{\RACER}{\textsc{Racer-GT}}
\newcommand{\indic}[1]{\mathbb{I}\!\left\{#1\right\}}
\newcommand{\Sig}{\bm{\Sigma}}
\newcommand{\wvec}{\bm{w}}
\newcommand{\evec}{\bm{e}}
\newcommand{\Zvec}{\bm{Z}}

% Design facets, used identically in both languages so that a reader can move
% between the Chinese and English documents without re-learning notation.
\newcommand{\facetT}{\mathcal{T}}   % historical date
\newcommand{\facetD}{\mathcal{D}}   % collection day
\newcommand{\facetS}{\mathcal{S}}   % collection stream
\newcommand{\facetR}{\mathcal{R}}   % technical replicate

\lstset{
  basicstyle=\ttfamily\footnotesize,
  breaklines=true,
  frame=single,
  rulecolor=\color{gray!40},
  backgroundcolor=\color{gray!5},
  columns=fullflexible,
  keepspaces=true,
  showstringspaces=false
}


\setlength{\parindent}{1.5em}

\newtheorem{assumption}{Assumption}[section]
\newtheorem{proposition}{Proposition}[section]
\newtheorem{corollary}{Corollary}[section]
\newtheorem{remark}{Remark}[section]
\newtheorem{definition}{Definition}[section]
\newtheorem{lemma}{Lemma}[section]

\lhead{RACER-GT: User Guide}

\title{\bfseries RACER-GT User Guide\\[0.4em]
\large Installation, Workflow, Data Contracts, and Outputs}
\author{\racerauthor\\\small \raceraffil}
\date{\racerdate\ \textbullet\ version \racerversion}

\begin{document}
\maketitle

\begin{abstract}
\noindent This guide takes a researcher from installation to an accepted daily
Google Trends index. It documents the data contract a collector must satisfy,
the four command-line steps, the Python API, every output file, and the failure
modes that are most likely to occur in practice. Method and proofs are in the
methodology manuscript and the mathematical appendix; this document is about
operating the software correctly.
\end{abstract}

\tableofcontents
\newpage

% =============================================================================
\section{Installation}

RACER-GT requires Python 3.10 or later.

\begin{lstlisting}[language=bash]
# from PyPI
python -m pip install racergt

# from a release wheel
python -m pip install racergt-1.0.0-py3-none-any.whl

# from source, for development
python -m pip install -e '.[dev]'
pytest -q
\end{lstlisting}

The package is deliberately \emph{ingestion-first}: it contains no Google Trends
scraper. Endpoints, rate limits, authentication, and terms of service change,
and a statistical method that depends on an unofficial client inherits that
fragility. The collector is your choice; the package only requires that it emit
the schema in \cref{sec:contract}.

% =============================================================================
\section{The four-step workflow}

\subsection{Step 1: lock the protocol}

\begin{lstlisting}[language=bash]
racergt design config.yaml --out protocol_bundle --anchor-date 2026-08-01
\end{lstlisting}

This produces \texttt{protocol.lock.yaml} with a SHA-256 protocol hash,
\texttt{collection\_schedule.csv} with the balanced day $\times$ stream $\times$
replicate plan, \texttt{chunk\_windows.csv} with the fixed query windows, and an
audit manifest.

\begin{remark}[Do not hand-edit the lock file]
\texttt{RacerGTConfig.load\_yaml} recomputes the hash and raises if it disagrees
with the stored value. Editing a locked setting by hand therefore fails loudly,
which is the intended behaviour. To change a setting, edit the source
configuration and regenerate.
\end{remark}

\subsection{Step 2: collect and convert}

Follow the schedule. Preserve every raw response unmodified, then convert to the
long format of \cref{sec:contract}. Record a hash of both the raw response and
the numeric vector for each chunk. If a network interruption, login change,
cookie clear, or client update occurs, record it in the metadata; do not decide
unilaterally that the pull is invalid, because that decision would be a
value-dependent exclusion.

\subsection{Step 3: audit}

\begin{lstlisting}[language=bash]
racergt audit config.yaml raw_chunks.csv --out audit.json
\end{lstlisting}

The audit checks the batch against the locked protocol before any estimation
happens. It exits with status 2 if the batch fails.

\subsection{Step 4: run}

\begin{lstlisting}[language=bash]
racergt run config.yaml raw_chunks.csv \
  --benchmark lower_frequency_benchmark.csv \
  --out results
\end{lstlisting}

Exit status is 3 if the locked decision tree returns FAIL, so the command can
gate an automated pipeline.

\subsection{Python API}

\begin{lstlisting}[language=Python]
import pandas as pd
from racergt import RacerGTConfig, RacerGTPipeline

config = RacerGTConfig.load_yaml("config.yaml")
raw = pd.read_csv("raw_chunks.csv")
benchmark = pd.read_csv("lower_frequency_benchmark.csv")

result = RacerGTPipeline(config).fit(raw, benchmark=benchmark)
result.save("results")

print(result.decision.status)      # PASS / REVIEW / FAIL
print(result.final_series.head())
\end{lstlisting}

% =============================================================================
\section{The data contract}\label{sec:contract}

\subsection{Raw chunk table}

One row per \texttt{pull\_id} $\times$ \texttt{chunk\_id} $\times$
\texttt{historical\_date}.

\begin{table}[H]
\centering\small
\caption{Required columns.}
\begin{tabularx}{\textwidth}{llX}
\toprule
Column & Type & Meaning\\
\midrule
\texttt{series\_id} & string & research indicator identifier\\
\texttt{pull\_id} & string & one complete history reconstruction\\
\texttt{chunk\_id} & string & fixed query window, must exist in the manifest\\
\texttt{historical\_date} & date & GT observation date\\
\texttt{value} & float & raw GT relative search value, normally 0--100\\
\texttt{collection\_day} & integer & design level of the collection day\\
\texttt{stream\_id} & string & fixed composite collection environment\\
\texttt{replicate\_id} & string & technical replicate within a day $\times$ stream cell\\
\texttt{window\_start} & date & chunk start\\
\texttt{window\_end} & date & chunk end\\
\bottomrule
\end{tabularx}
\end{table}

\begin{remark}[Two columns that are routinely misunderstood]
\texttt{collection\_day} is the \emph{ordinal design level} of the collection
day (0, 1, 2, \dots), not the day offset and not a calendar date. And
\texttt{stream\_id} denotes a composite environment---device, browser profile,
cookies, account state, network route, IP. An estimated stream effect must not
be reported as an IP effect unless a crossover factorial design was actually
run.
\end{remark}

Strongly recommended additional columns: \texttt{retrieved\_at},
\texttt{protocol\_hash}, \texttt{keyword}, \texttt{geo}, \texttt{category},
\texttt{search\_property}, \texttt{language}, \texttt{is\_partial},
\texttt{raw\_response\_hash}, \texttt{network\_event}.

\subsection{Lower-frequency benchmark table}

Requires \texttt{series\_id}, \texttt{period\_start}, \texttt{period\_end},
\texttt{value}; \texttt{se} is optional and falls back to the protocol default.

% =============================================================================
\section{Outputs}

\begin{table}[H]
\centering\small
\caption{Files written by \texttt{result.save()}.}
\begin{tabularx}{\textwidth}{lX}
\toprule
Path & Contents\\
\midrule
\texttt{RACER\_GT\_final\_series.csv} & final daily index, conditional SE, 95\% interval\\
\texttt{complete\_calibrated\_pulls.csv} & every calibrated pull\\
\texttt{calibration/*\_chunk\_scales.csv} & per-chunk global scale estimates\\
\texttt{calibration/*\_overlap\_edges.csv} & robust edge estimates and diagnostics\\
\texttt{duplicates/} & exact groups, pair metrics, components, residuals\\
\texttt{consensus/design\_weighted\_consensus.csv} & design-based reference estimator\\
\texttt{consensus/consensus\_weights.csv} & GLS / minimum-variance weights\\
\texttt{gstudy/} & ANOVA, variance components, D-study, per transformation\\
\texttt{reliability/} & pairwise reliability, convergence, day/stream dependence\\
\texttt{decision/acceptance\_decision\_tree.csv} & every rule, observed vs threshold\\
\texttt{report/RACER\_GT\_diagnostic\_report.md} & figure-based diagnostic report\\
\texttt{summary.json} & protocol hash, status, headline diagnostics\\
\bottomrule
\end{tabularx}
\end{table}

\begin{remark}[Two consensus files, one final series]
\texttt{consensus/design\_weighted\_consensus.csv} is a design-based reference
estimator. It is reported for comparison and does \emph{not} feed
\texttt{RACER\_GT\_final\_series.csv}, which comes from the temporal benchmark
when one is supplied and from the GLS consensus otherwise. Large disagreement
between the two is a warning sign worth investigating, not a bug.
\end{remark}

% =============================================================================
\section{Reading the decision}

A batch is FAIL if any mandatory rule fails, REVIEW if none fails but at least
one is indeterminate, and PASS otherwise. Common causes, and what to do:

\begin{table}[H]
\centering\small
\caption{Frequent outcomes and their remedies.}
\begin{tabularx}{\textwidth}{lX}
\toprule
Symptom & Interpretation and action\\
\midrule
\texttt{OverlapGraphError} & the chunk design leaves a gap. Increase window length or reduce step so that consecutive chunks share at least \texttt{min\_overlap\_days} usable days.\\
Low spectral effective pulls & retrievals are highly dependent. Add collection days or streams; adding technical replicates will not help.\\
Detection rule indeterminate & no zero/non-zero variation, so $\kappa$ is undefined. The rule is automatically non-mandatory in this case and needs no action.\\
High zero share & the term is too rare for daily frequency. Move to weekly, or use a broader Topic rather than a Term.\\
Innovation G below threshold, level G above & levels reproduce but daily changes do not. Do not use the series in differences.\\
Benchmark RMSE too large & the daily consensus and the lower-frequency series disagree. Check that both use the identical query specification.\\
\bottomrule
\end{tabularx}
\end{table}

% =============================================================================
\section{Stata interoperability}

Every table can be exported to Stata 118 format:

\begin{lstlisting}[language=bash]
racergt export-stata results/RACER_GT_final_series.csv final_series.dta
\end{lstlisting}

Supported read formats are CSV, Parquet, and XLS/XLSX; supported write formats
are CSV, Parquet, XLSX, and Stata 118 \texttt{.dta}.

% =============================================================================
\section{Research conduct requirements}

These are not optional style preferences; violating them invalidates the
guarantees the framework provides.

\begin{itemize}
\item A static IP address is not evidence of an independent sample.
\item Do not adjust reliability thresholds or select pulls after seeing
financial results.
\item If pilot monitoring leads to any protocol or decision-rule change, the
monitored data become pilot data and the confirmatory batch restarts.
\item Retain exact duplicates in the audit archive.
\item PASS means the batch satisfies the locked measurement rules. It does not
establish construct validity for the keyword and licenses no causal claim.
\end{itemize}

\end{document}
